% !TEX root = ../main.tex
\begin{apendicesenv}
\begin{otherlanguage*}{brazil}

\chapter{Declaração de Uso de Ferramentas de Inteligência Artificial}
\label{apendice:declaracao_ia}

Em conformidade com as diretrizes de integridade acadêmica e transparência científica da Universidade de Brasília (UnB), declara-se que foram utilizadas ferramentas de Inteligência Artificial (IA) generativa como suporte técnico-redacional durante a elaboração deste Trabalho de Conclusão de Curso:

\begin{itemize}
    \item \textbf{Perplexity AI}: utilizada como instrumento de consulta e checagem factual para a verificação das informações técnicas, conceituais e bibliográficas escritas no documento;
    \item \textbf{Claude (Anthropic)}: utilizado para auxílio na revisão estilística e correção de erros gramaticais e ortográficos da língua inglesa ao longo do manuscrito.
\end{itemize}

Ressalta-se que o uso das ferramentas se ateve exclusivamente ao aprimoramento redacional e à verificação das informações redigidas. Toda a fundamentação teórica, o desenvolvimento do código-fonte do simulador, a concepção da metodologia experimental, a execução dos ensaios computacionais e a análise crítica dos resultados obtidos foram integralmente realizados pelo autor.

O autor assume plena e exclusiva responsabilidade pelo conteúdo, exatidão, veracidade e integridade intelectual de todo o trabalho apresentado, atestando que nenhuma ferramenta de inteligência artificial figura como autora ou coautora deste documento.

\vspace{1.5cm}

\begin{center}
    \imprimirlocal, \imprimirdata.

    \vspace{1.5cm}

    \begin{minipage}{0.65\textwidth}
        \centering
        \rule{\textwidth}{0.5pt} \\
        \textbf{\imprimirautor} \\
        Discente do Curso de Engenharia de Software \\
        Faculdade de Ciências e Tecnologias em Engenharia (FCTE) \\
        Universidade de Brasília (UnB)
    \end{minipage}
\end{center}

\end{otherlanguage*}
\end{apendicesenv}
